\section{Conclusion}

We set out to answer: when AI assistants match user communication style, do users engage more? The answer is nuanced. AI systems do accommodate ($r=0.152$), and accommodation relates to engagement---but the relationship is non-linear.

\textbf{Moderate accommodation predicts highest continuation rates (71.4\%)}, while both high (65.9\%) and low (60.9\%) accommodation underperform. This ``Goldilocks zone'' suggests AI should adapt ``just right''---neither too little nor too much.

The implications are practical: chatbot designers should calibrate accommodation rather than maximize it. The implications are theoretical: CAT predictions require qualification in human-AI contexts where one party is perceived as artificial.

The Goldilocks zone is not a limitation---it is design guidance.
